\chapter{Конструкторская часть}

\section{Требования к программному обеспечению}

К разрабатываемому программному обеспечению предъявляются следующие требования:

\begin{enumerate}
    \item ввод линейных размеров матриц $M, N, Q$ (натуральные числа);
    \item ввод целочисленных элементов матриц $A$ и $B$ вручную с клавиатуры либо автоматическая случайная генерация элементов в заданном числовом диапазоне;
    \item проверка корректности размерностей матриц (число столбцов матрицы $A$ должно совпадать с числом строк матрицы $B$);
    \item вывод текущих матриц $A$ и $B$ в консоль в форматированном виде;
    \item циклический вывод интерактивного меню с возможностью выбора и выполнения следующих операций:
    \begin{enumerate}[label=\arabic{enumi}.\arabic*.]
        \item вычисление произведения матриц с использованием стандартного алгоритма;
        \item вычисление произведения матриц с использованием базового алгоритма Винограда;
        \item вычисление произведения матриц с использованием оптимизированного по варианту алгоритма Винограда;
        \item перегенерация исходных матриц случайными значениями;
        \item повторный ввод размерностей и элементов матриц пользователем;
        \item вывод текущих матриц на экран;
        \item запуск серии массированных замеров процессорного времени выполнения всех трёх алгоритмов для квадратных матриц чётных и нечётных размеров с выгрузкой результатов в CSV-файлы;
        \item завершение работы программы.
    \end{enumerate}
\end{enumerate}

\section{Разработка алгоритмов}

Для каждого из разрабатываемых алгоритмов умножения определены:
\begin{itemize}
    \item $A$ --- входная матрица размера $M \times N$ с целочисленными элементами;
    \item $B$ --- входная матрица размера $N \times Q$ с целочисленными элементами;
    \item $C$ --- результирующая матрица размера $M \times Q$, элементы которой перед началом вычислений инициализируются нулями;
    \item $M, N, Q$ --- линейные размеры перемножаемых матриц.
\end{itemize}

Для базового и оптимизированного алгоритмов Винограда дополнительно определены:
\begin{itemize}
    \item $row$ --- вспомогательный вектор длины $M$ для промежуточных сумм парных произведений строк матрицы $A$;
    \item $col$ --- вспомогательный вектор длины $Q$ для промежуточных сумм парных произведений столбцов матрицы $B$;
    \item $half\_n$ (в оптимизированной версии) --- кэшированная граница диапазонов циклов $\lfloor N / 2 \rfloor$;
    \item $two\_k$ (в оптимизированной версии) --- локальная переменная, хранящая вычисленный сдвигом индекс $k \ll 1$ ($2k$).
\end{itemize}

\subsection{Стандартный алгоритм умножения матриц}

Стандартный алгоритм производит прямое вычисление каждого элемента $C_{i,j}$ результирующей матрицы по математическому определению произведения матриц[cite: 103, 143]:
\begin{equation}
    C_{i,j} = \sum_{k=0}^{N-1} A_{i,k} \cdot B_{k,j}, \quad i = 0, \ldots, M-1, \quad j = 0, \ldots, Q-1.
\end{equation}

Алгоритм состоит из трёх вложенных циклов[cite: 103, 143]: внешний цикл организует проход по строкам $i \in [0, M - 1]$, средний --- по столбцам $j \in [0, Q - 1]$, а внутренний выполняет накопление суммы произведений по индексу $k \in [0, N - 1]$[cite: 104, 105, 147]. Схема стандартного алгоритма умножения матриц приведена на рисунке~\ref{img:algo-standard}[cite: 103, 146].

\begin{figure}[H]
    \centering
    \includegraphics[scale=0.85]{image/algo/algo_1.pdf}
    \caption{Схема стандартного алгоритма умножения матриц}
    \label{img:algo-standard}
\end{figure}

\subsection{Алгоритм Винограда}

Алгоритм Винограда уменьшает общее число скалярных умножений элементов матриц ценой увеличения количества операций сложения[cite: 105, 144]. Вычисление произведения основано на предварительном расчёте двух одномерных массивов, зависящих только от строк матрицы $A$ и столбцов матрицы $B$ соответственно[cite: 106, 144]:
\begin{equation}
    row_i = \sum_{k=0}^{\lfloor N/2 \rfloor - 1} A_{i, 2k} \cdot A_{i, 2k+1}, \quad i = 0, \ldots, M - 1,
\end{equation}
\begin{equation}
    col_j = \sum_{k=0}^{\lfloor N/2 \rfloor - 1} B_{2k, j} \cdot B_{2k+1, j}, \quad j = 0, \ldots, Q - 1.
\end{equation}

После вычисления массивов $row$ и $col$ каждый элемент матрицы результата инициализируется значением $-row_i - col_j$, после чего в основном цикле суммируются произведения парных сумм элементов[cite: 106, 144]:
\begin{equation}
    C_{i,j} = -row_i - col_j + \sum_{k=0}^{\lfloor N/2 \rfloor - 1} (A_{i, 2k} + B_{2k+1, j}) \cdot (A_{i, 2k+1} + B_{2k, j}).
\end{equation}

Если общий размер $N$ является нечётным числом, последний непарный элемент строки и столбца учитывается отдельной коррекцией[cite: 106, 144]:
\begin{equation}
    C_{i,j} = C_{i,j} + A_{i, N-1} \cdot B_{N-1, j}.
\end{equation}

Схема базового алгоритма Винограда приведена на рисунке~\ref{img:algo-winograd}[cite: 106, 148].

\begin{figure}[H]
    \centering
    \includegraphics[scale=0.85]{image/algo/algo_2.pdf}
    \caption{Схема базового алгоритма Винограда}
    \label{img:algo-winograd}
\end{figure}

\subsection{Оптимизированный алгоритм Винограда}

В соответствии с индивидуальным вариантом к алгоритму Винограда применены следующие две оптимизации[cite: 112, 193]:
\begin{enumerate}[label=\alph*)]
    \item замена умножения $2 \cdot k$ при индексации побитовым сдвигом влево на один разряд: $k \ll 1$[cite: 112, 193]. Значение вычисляется один раз в начале каждой итерации цикла по $k$, сохраняется в локальную переменную $two\_k$ и затем повторно используется во всех обращениях к элементам текущей итерации[cite: 112, 134, 135];
    \item предварительное вычисление (кэширование) значения верхней границы циклов $half\_n = \lfloor N / 2 \rfloor$ в отдельную константу до входа в циклы вместо многократного повторения целочисленного деления $N / 2$ при каждой проверке условия продолжения цикла[cite: 112, 134, 193].
\end{enumerate}

Применение данных оптимизаций не изменяет результирующие значения матрицы $C$ и общее количество операций над её элементами, однако сокращает накладные расходы на целочисленную арифметику внутри циклов[cite: 112]. Схема оптимизированного алгоритма Винограда приведена на рисунке~\ref{img:algo-winograd-opt}[cite: 112].

\begin{figure}[H]
    \centering
    \includegraphics[scale=0.85]{image/algo/algo_3.pdf}
    \caption{Схема оптимизированного алгоритма Винограда}
    \label{img:algo-winograd-opt}
\end{figure}

\section{Вывод}

В результате выполнения конструкторской части были формализованы требования к разрабатываемому программному обеспечению, описаны математические основы и спецификации входных и выходных данных для стандартного алгоритма умножения матриц, базового алгоритма Винограда и его оптимизированной модификации, а также представлены схемы данных алгоритмов[cite: 12].